\section{Matrix Landscape: Discovery, Adapters, and Scope}
\label{sec:results}

Table~\ref{tab:main-matrix} is a workload landscape, not an efficiency or selector-attribution
result. Under nominal budget $B{=}6$, it records Class~A generators and Class~B identical-pool
selectors across providers and datasets; Section~\ref{sec:setup} defines the unscored
Fireworks$\times$GPQA protocol-blocked outcome. Sections~\ref{sec:identical-pool}--\ref{sec:compute}
then isolate the comparisons the matrix alone cannot settle: identical-pool selection, provider
transfer, and realized-resource accounting. Bounded repair follows the historical convention in
Section~\ref{sec:method-scope}; \ref{app:repair} verifies that a uniform-repair
counterfactual causes no external winner flips.

\begin{table}[tbp]
\centering
\footnotesize
\setlength{\tabcolsep}{3pt}
\caption{Provider$\times$dataset landscape under nominal $B{=}6$. Class~A columns are closed-API
generators (Frontier diagnostic; L1-/S1-/TALE-style adapters). Class~B columns select over the
\emph{same} four-generator pool (pool cost $\leq 4B$; $0$ extra model calls). Bold marks the highest
Class~A count and, separately, the higher of FTA vs.\ Pooled-4 as descriptive highlighting. Trust
tiers and the Fireworks$\times$GPQA blocked outcome are defined in
Section~\ref{sec:setup}.}
\label{tab:main-matrix}
% Class A: Frontier plus L1/S1/TALE adapters (nominal B=6 each).
% Class B: FTA replay and Pooled-4 (same shared pool; zero marginal selection calls).
% Bold = max Class A count; separately, higher of FTA vs Pooled-4 (both if tie).
\begin{tabular}{@{}llrrrrrr@{}}
\toprule
Provider & Data & Front. & L1 & S1 & TALE & FTA & Pooled-4 \\
\midrule
Cohere & GSM8K & 230/300 & \textbf{249/300} & 246/300 & 235/300 & 260/300 & \textbf{261/300} \\
Cohere & MATH-500 & \textbf{87/300} & 73/300 & 84/300 & 76/300 & 92/300 & \textbf{98/300} \\
Cohere & GPQA-D & 68/198 & 74/198 & 69/198 & \textbf{77/198} & \textbf{72/198} & \textbf{72/198} \\
Cohere & StrategyQA & 78/100 & 69/100 & 77/100 & \textbf{83/100} & \textbf{76/100} & \textbf{76/100} \\
Azure & GSM8K & \textbf{276/300} & 269/300 & 246/300 & 204/300 & 260/300 & \textbf{276/300} \\
Azure & MATH-500 & \textbf{167/300} & 142/300 & 161/300 & 112/300 & 143/300 & \textbf{164/300} \\
Azure & GPQA-D & 103/198 & \textbf{105/198} & 104/198 & 94/198 & 102/198 & \textbf{104/198} \\
Azure & StrategyQA & 70/100 & 72/100 & \textbf{76/100} & 74/100 & \textbf{72/100} & \textbf{72/100} \\
Vertex & GSM8K & \textbf{279/300} & 262/300 & 244/300 & 255/300 & 272/300 & \textbf{278/300} \\
Vertex & MATH-500 & 217/300 & 220/300 & 210/300 & \textbf{223/300} & \textbf{228/300} & 224/300 \\
Vertex & GPQA-D & \textbf{113/198} & 106/198 & 108/198 & 110/198 & \textbf{125/198} & 123/198 \\
Vertex & StrategyQA & \textbf{84/100} & 82/100 & 83/100 & \textbf{84/100} & \textbf{84/100} & \textbf{84/100} \\
Fireworks & GSM8K & \textbf{223/300} & 210/300 & 160/300 & 191/300 & 224/300 & \textbf{230/300} \\
Fireworks & MATH-500 & 100/300 & 105/300 & 94/300 & \textbf{114/300} & \textbf{117/300} & \textbf{117/300} \\
Fireworks & GPQA-D & \multicolumn{6}{c}{BLOCKED\_PROTOCOL\_NONCONVERGENCE} \\
Fireworks & StrategyQA & \textbf{81/100} & 80/100 & 80/100 & 80/100 & \textbf{79/100} & \textbf{79/100} \\
\bottomrule
\end{tabular}

\end{table}

\subsection{Does higher discovery predict safer selection?}

Absent-from-pool errors are unrecoverable by pool-only selectors; present-but-misselected errors are
the only errors a commit rule can still fix (Equation~\eqref{eq:acc-decomp}). Writing $\mathrm{DR}$
for discovery rate and $\mathrm{CSA}$ for conditional selection accuracy
(Equation~\eqref{eq:rates}; Figure~\ref{fig:protocol}), higher discovery does not systematically track
FTA gains: global Pearson $r=-0.185$ (not significant at $n{=}15$; unstable under leave-one-out;
supplementary scatter diagnostic). Azure$\times$GSM8K has near-ceiling discovery with a large negative FTA
delta, while Vertex and Azure GPQA share essentially the same discovery rate with opposite selector
signs (Section~\ref{sec:provider-analysis}). Joint $\mathrm{DR}$/$\mathrm{CSA}$ reporting is
therefore required before identical-pool tests.

\subsection{What do Class~A adapter counts characterize?}

Descriptively, within the evaluated closed-API adapter surfaces, the reserved-attempt full-expansion
diagnostic is sole highest in $7/15$ cells and tied highest in one more; remaining adapter surfaces
occupy fewer cells. We use these counts to characterize the workload surface before the controlled
selection and accounting analyses. An offline incumbent-only counterfactual loses to full Frontier
in $9/15$ cells (Section~\ref{sec:ablations}), indicating that challenger expansion is not inert
under stored traces.

\subsection{What does the single-cell same-model control establish?}
\label{sec:sc-entitlement-control}

Table~\ref{tab:compute-matched-phase7} is a single-cell supporting control on Azure GPT-4.1-mini
with GSM8K ($n{=}300$): it matches call/sample budget between Frontier
($B{=}6$) and true same-model fixed-$N$ self-consistency (SC; six samples). Each obtains $276/300$
($92.00\%$; marginal exact binomial 95\% CI $88.3304$--$94.8071\%$). Paired records show 298 ties
and one method-only correct case each (two-sided exact McNemar $p{=}1.0$; paired-bootstrap 95\%
percentile interval for Frontier$-$SC: $[-1.00,+1.00]$\,pp). No equivalence margin was prespecified.
Panel~B is contextual only.

\begin{table}[tbp]
\centering
\footnotesize
\caption{Single-cell supporting control on Azure \texttt{gpt-4.1-mini}$\times$GSM8K ($n{=}300$).
Panel~A matches call/sample budget (Frontier $B{=}6$ vs.\ SC with six samples); rows are not
token-, latency-, dollar-, or FLOP-matched. Panel~B is contextual adapter/FTA context on the same
IDs.}
\label{tab:compute-matched-phase7}
% Values are unchanged from the corrected common-set method summary.
% Panel A is the direct paired control. Panel B is contextual and mixes generation adapters
% with an offline pool selector; the operational columns prevent a false equal-cost reading.
\begin{tabularx}{0.98\textwidth}{@{}p{2.35cm}Yp{1.15cm}Yrrr@{}}
\toprule
Method & Operational class & Calls? & Pool/call basis & Correct & Total & Accuracy \\
\midrule
\multicolumn{7}{@{}l}{\textit{Panel A: direct paired same-model control (common 300 examples)}} \\
Frontier & Structured generation & Yes & Nominal budget $B{=}6$ & 276 & 300 & 92.00\% \\
True same-model SC $N{=}6$ & Fixed-$N$ generation + vote & Yes & Exactly $6$ samples & 276 & 300 & 92.00\% \\
\midrule
\multicolumn{7}{@{}l}{\textit{Panel B: contextual adapter and replay rows (not a cost-matched class)}} \\
L1-style adapter & Generation adapter & Yes & Nominal budget $B{=}6$ & 269 & 300 & 89.67\% \\
FTA replay selector & Offline pool selector & No & Shared four-generator pool ($\leq4B$); no selector calls & 260 & 300 & 86.67\% \\
S1-inspired adapter & Generation adapter & Yes & Nominal budget $B{=}6$ & 246 & 300 & 82.00\% \\
TALE-EP-style adapter & Generation adapter & Yes & Nominal budget $B{=}6$ & 204 & 300 & 68.00\% \\
\bottomrule
\end{tabularx}

\end{table}
